\documentclass[12pt]{article}
\usepackage[utf8]{inputenc}
\usepackage[T1]{fontenc}
\usepackage[margin=1in]{geometry}
\usepackage{setspace}
\usepackage{xcolor}
\usepackage[colorlinks=true, urlcolor=blue!60!black, linkcolor=blue!60!black]{hyperref}
\setlength{\parindent}{2em}
\setlength{\parskip}{1.25em}
\renewcommand{\baselinestretch}{1.0}
\begin{document}
\begin{center}
{\LARGE \textit{FDA v. Alliance for Hippocratic Medicine}, 602 U.S. 367 (2024)\par}
\vspace{0.8em}
{\large Case Brief --- Four-Page Edition\par}
\end{center}
\vspace{1em}

\section{Memory Jogger}

Pro-life doctors and medical associations who neither prescribe nor use mifepristone lack standing to challenge the FDA's 2016 and 2021 decisions relaxing access to it: conscience laws break the causal chain, there is no ``doctor standing,'' and an organization ``cannot spend its way into standing'' (\href{https://www.supremecourt.gov/opinions/23pdf/602us1r35_bqmd.pdf}{602 U.S. 367}).

\section{Detailed Case Facts}

The FDA approved mifepristone (Mifeprex) in 2000 for terminating pregnancies up to seven weeks, with strict conditions: physician prescribing, three in-person visits, and reporting of serious adverse events. In 2016 it extended the limit to ten weeks, allowed non-physician prescribers, reduced in-person visits to one, and limited reporting to fatalities; in 2019 it approved a generic; and in 2021 it stopped enforcing in-person dispensing, permitting telemedicine and mail. In November 2022, after \textit{Dobbs}, four pro-life medical associations and several doctors sued under the APA. None of them prescribed, sold, or used mifepristone, and the FDA required nothing of them. Danco Laboratories intervened to defend the FDA.

\section{Procedural History}

Judge Kacsmaryk of the Northern District of Texas found standing and stayed the 2000 approval and all later actions, effectively taking mifepristone off the market. The Fifth Circuit narrowed the stay, and the Supreme Court stayed the order entirely pending review. On the merits the Fifth Circuit found standing and upheld the injunction as to the 2016 and 2021 actions. The Supreme Court granted certiorari and reversed on June 13, 2024.

\section{Judicial Votes}

Reversed and remanded, \textbf{9--0}.
\begin{itemize}
\item \textbf{Opinion of the Court:} Justice Kavanaugh, for a unanimous Court.
\item \textbf{Concurring:} Justice Thomas, questioning associational standing.
\item \textbf{Dissents:} None.
\end{itemize}

\section{Holding}

Unregulated plaintiffs who do not prescribe, make, sell, or use a drug lack standing to challenge the FDA's relaxation of restrictions on it. (1) Federal conscience laws, which EMTALA does not override, break any causal chain to conscience injuries. (2) Downstream costs of treating others' complications are too speculative and attenuated; there is no ``doctor standing.'' (3) Under a narrowed \textit{Havens Realty}, an organization cannot create standing by spending money to study or oppose a policy that does not impede its core activities. (4) That no one else might have standing does not create standing. The merits were not reached.

\section{Analysis of Opinions}

\textbf{Majority (Kavanaugh).} Invoking Justice Scalia's question, ``What's it to you?,'' Kavanaugh stressed that courts do not hear general policy objections. When an unregulated plaintiff challenges the regulation of others, causation is substantially more difficult to establish, and the chain may be neither speculative nor attenuated. \emph{Conscience:} being forced to act against conscience is an injury, but the Church Amendments protect doctors, and no plaintiff identified a single instance of compulsion since 2000; mere distress at others' conduct is not enough. \emph{Economic harm:} claims of diverted time and higher malpractice or insurance costs lacked support, and recognizing doctor standing would let physicians, or firefighters and teachers, challenge any safety deregulation. \emph{Organizations:} \textit{Havens} involved direct interference with an organization's core business, like a retailer sold defective goods; it is an unusual case not to be extended. Some questions are left to the political branches.

\textbf{Thomas (concurring).} The Court should reconsider associational standing, which lets an uninjured association sue for a single member---here an association of associations---and encourages universal injunctions, evades Rule 23, and arose almost by accident in 1963.

\section{Examples: Future Applications}

\begin{enumerate}
\item \textbf{Same-Side (No Standing): Voter-Roll Rules.} Advocacy groups divert resources to retrain volunteers against a voter-registration law. They cannot spend their way into standing, as the Ninth Circuit held in \textit{Arizona Alliance for Retired Americans v. Mayes}.
\item \textbf{Same-Side (No Standing): Scooter Deregulation.} ER doctors challenge repeal of scooter safety standards, predicting more injured riders. No doctor standing.
\item \textbf{Opposite-Side (Standing): Rival Manufacturer.} A competitor challenges an unlawful approval of a cheaper rival drug---a preserved category of competitor standing.
\item \textbf{Opposite-Side (Standing): Directly Regulated Doctor.} A state threatens license revocation for refusing a procedure; the doctor is directly regulated.
\item \textbf{Fence-Sitter: Legal-Aid Nonprofit.} A group providing direct representation is barred from detained clients or fed false case information. It can analogize to \textit{Havens}, but courts ask whether the policy obstructs core operations or merely makes them costlier.
\end{enumerate}

\section{Critique}

\textbf{Originalist views.} Conservatives praise a disciplined restatement of Article III and the narrowing of resource-diversion and offended-observer standing; Thomas would go further. Observers note that the Solicitor General's concessions on conscience protections mooted that theory while cementing strong conscience rights.

\textbf{Progressive views.} The \textit{Harvard Law Review} notes that the decision left the merits open to suits by states such as Missouri, Kansas, and Idaho, especially after \textit{Loper Bright} and \textit{Corner Post}. Critics contrast its rigor with looser standing in \textit{Biden v. Nebraska} and \textit{303 Creative}, and warn that the narrowing of \textit{Havens} will curtail civil-rights and environmental nonprofits. Strict traceability can also leave deregulation unreviewable.

\section{Key Quotations}

\begin{quote}
``But an organization that has not suffered a concrete injury caused by a defendant's action cannot spend its way into standing simply by expending money to gather information and advocate against the defendant's action. An organization cannot manufacture its own standing in that way.''
\end{quote}
\noindent --- Justice Kavanaugh, for the Court, 602 U.S. at 394.

\begin{quote}
``The injury in fact requirement prevents the federal courts from becoming a `vehicle for the vindication of the value interests of concerned bystanders.' \textit{Allen}, 468 U. S., at 756 (quotation marks omitted). An Article III court is not a legislative assembly, a town square, or a faculty lounge. Article III does not contemplate a system where 330 million citizens can come to federal court whenever they believe that the government is acting contrary to the Constitution or other federal law.''
\end{quote}
\noindent --- Justice Kavanaugh, for the Court, 602 U.S. at 382.

\begin{quote}
``Stated otherwise, there is no Article III doctrine of `doctor standing' that allows doctors to challenge general government safety regulations. Nor will this Court now create such a novel standing doctrine out of whole cloth.''
\end{quote}
\noindent --- Justice Kavanaugh, for the Court, 602 U.S. at 391.

\end{document}
